\section{Sharp space with higher-dimensional density-matrix outputs}
\label{sec:matrixspace65}
Positive rounding removes the restriction of the uniform upper
construction to a Bloch ball. We retain the spectral hypotheses in the
converse; none is needed by the algorithm.

\begin{theorem}[Sharp matrix-output space]\label{thm:matrixspace65}
Fix a finite family of Gaussian-rational unitary matrices in $U(d)$
and a Gaussian-rational pure initial state $\rho_0$. Let $G$ be the
compact closure of their conjugation action on the real Euclidean
space
\[
 E=\{H\in\C^{d\times d}:H=H^*,\ \tr H=0\},
 \qquad \|H\|=\|H\|_{\rm F}.
\]
Assume a probability law supported on the given commands has a full
Koopman action gap on the unit sphere of $E$, off
\emph{all} invariant functions, and an all-direction cap estimate
\eqref{eq:caps62} of exponent $p>0$. The dimension, matrices, seed,
rational denominators, cap constants and spectral constant are fixed
independently of $N,L$.

For a single fixed-program classical one-pass simulator, assume
almost-sure termination between commands and at final output, and require a
legal numerical output $\widehat\rho(w)\in\mathcal D_d$ on every run
and only the mean condition
\begin{equation}\label{eq:matrixmean65}
 \sup_{|w|=N}\|\E\widehat\rho(w)-U_w\rho_0U_w^*\|_{\rm F}\le2^{-L}.
\end{equation}
For every such fixed program, a matching space lower bound holds for
$L\ge2$ and all sufficiently large $N$, with constants and threshold
allowed to depend on that program and the fixed physical data.
A single uniform algorithm attains the upper bound. In this
fixed-program asymptotic sense, the optimal writable bit-space has order
\begin{equation}\label{eq:matrixspace65}
          \Theta_{\rm input}(\min\{N,L+\log_2(N+1)\}).
\end{equation}
The lower bound permits arbitrary randomization and unlimited
internal running time. A deterministic integer program attains the
upper bound with pointwise trace-norm error. Exact computation has
order $N$. No idle or return assumption is made.
\end{theorem}
\begin{proof}
Put $c=I/d$ and $r_d=\sqrt{(d-1)/d}$. For every density matrix $D$,
\[
 \|D-c\|_{\rm F}^2=\tr(D^2)-1/d\le r_d^2;
\]
equality holds for a pure state. Thus the physical target
$X_w=(U_w\rho_0U_w^*-c)/r_d$ is unit, and every normalized legal
decoder $(D-c)/r_d$ has norm at most one.

The complete-configuration reduction of
Lemma~\ref{lem:boundary64} works with $\mathcal D_d$ in place of
$\mathcal D_2$: almost-sure termination gives stochastic rows, and
conditional means of legal outputs remain legal by compact convexity.
The induced rows may depend on the public boundary index, which
only enlarges the lower-bound comparison class.
The proof of Theorem~\ref{thm:driftoccupation62} uses the decoder
set only through the unit-norm bound. With error
$\varepsilon/r_d$ and $0\le\varepsilon\le1/4$, its finite budget
therefore gives
\begin{equation}\label{eq:matrixbudget65}
 N\le C_0+C_1\log k+C_2\varepsilon k^{2/p}
\end{equation}
for a simulator with at most $k$ complete private configurations at
every positive boundary. Constants depend only on the fixed data.
Indeed the residual mass is $1-\varepsilon/r_d$, bounded away from
zero uniformly in this error range.

For $N\ge2C_0$, either
$\log k\ge N/(4C_1)$ or
$k^{2/p}\ge N/(4C_2\varepsilon)$.
Taking logarithms yields
\[
 \log_2k\ge c\min\{N,L+\log_2(N+1)\}-C'.
\]
Since the minimum tends to infinity with $N$, uniformly for $L\ge2$,
the additive constant is absorbed above a fixed physical-data-dependent horizon.
For a fixed finite program, the complete-configuration count is at
most $C(s+1)^a2^{bs}$ when $s$ work bits are used, with fixed
machine-dependent constants. After absorbing the fixed finite-control factor $C$, this proves the
work-space lower bound above a program-dependent horizon.
No bound uniform over arbitrarily large uncharged finite programs is
asserted.
At zero error \eqref{eq:matrixbudget65} gives $\log k\ge cN-C'$.

Regard each unitary as a one-outcome rational Kraus instrument.
Corollary~\ref{cor:channelstream65} supplies the deterministic
trace-norm upper bound and the exact fallback. Trace norm dominates
Frobenius norm. The algorithm does not use the spectral or cap
constants, nor construct a stochastic orbit-hull table.
\end{proof}

\begin{corollary}[An explicit family in every fixed dimension]
\label{cor:alldim65}
For every fixed $d\ge2$, take the adjacent-coordinate rational
unitaries displayed in Section~\ref{sec:matrix57}, their inverses and
the identity, and take $\rho_0=e_1e_1^*$. Legal numerical density
matrix simulation satisfies \eqref{eq:matrixspace65}, with
dimension-dependent constants.
\end{corollary}
\begin{proof}
The density argument in that section shows that the generated compact
group is $SU(d)$. All entries are algebraic. The
Bourgain--Gamburd theorem \cite{BG2012}, alternatively the
Benoist--de Saxc\'e extension \cite[Theorem~1.2]{BdS65}, gives a
regular-representation spectral gap for this fixed dense algebraic
set. Symmetrization and laziness give a contraction norm strictly
less than one. The identity is already an allowed command, so this
is an averaging measure on actual letters. Peter--Weyl decomposition
transfers that norm to the action off its invariant functions.
For conjugation on traceless Hermitian matrices, the fixed space
is zero. The least orbit dimension calculation in
Theorem~\ref{thm:minorbit58} gives $p_*=2(d-1)$ and its uniform
cap estimate. Theorem~\ref{thm:matrixspace65} applies.
\end{proof}
This corollary does not supply a numerical spectral constant uniformly
in $d$. It uses an explicit rational alphabet, not an unspecified
existence argument. The construction for a given alphabet is effective
without knowing its spectral gap; that gap is a premise of the lower
proof only.

\begin{corollary}[Sharpness for an instrument family containing an expanding unitary subsystem]
\label{cor:instrumentsharp65}
Suppose the fixed rational instrument family of
Theorem~\ref{thm:instrumentstream65} contains the unitary commands
and pure seed of Theorem~\ref{thm:matrixspace65}. Among simulators
that additionally return a legal numerical matrix and satisfy the
classical--quantum error requirement \eqref{eq:instrumenterror65},
the optimal space has order \eqref{eq:matrixspace65}.
\end{corollary}
\begin{proof}
The instrument simulator gives the upper bound. For the lower bound,
restrict the controller to a deterministic word of unitary commands.
Their outcome sets are singletons, so the error requirement implies
\eqref{eq:matrixmean65}. Apply the preceding theorem. This is not a
lower bound for transcript-only simulators without a legal numerical
matrix interface.
\end{proof}
